% Options for packages loaded elsewhere
\PassOptionsToPackage{unicode}{hyperref}
\PassOptionsToPackage{hyphens}{url}
%
\documentclass[
]{article}
\usepackage{amsmath,amssymb}
\usepackage{iftex}
\ifPDFTeX
  \usepackage[T1]{fontenc}
  \usepackage[utf8]{inputenc}
  \usepackage{textcomp} % provide euro and other symbols
\else % if luatex or xetex
  \usepackage{unicode-math} % this also loads fontspec
  \defaultfontfeatures{Scale=MatchLowercase}
  \defaultfontfeatures[\rmfamily]{Ligatures=TeX,Scale=1}
\fi
\usepackage{lmodern}
\ifPDFTeX\else
  % xetex/luatex font selection
\fi
% Use upquote if available, for straight quotes in verbatim environments
\IfFileExists{upquote.sty}{\usepackage{upquote}}{}
\IfFileExists{microtype.sty}{% use microtype if available
  \usepackage[]{microtype}
  \UseMicrotypeSet[protrusion]{basicmath} % disable protrusion for tt fonts
}{}
\makeatletter
\@ifundefined{KOMAClassName}{% if non-KOMA class
  \IfFileExists{parskip.sty}{%
    \usepackage{parskip}
  }{% else
    \setlength{\parindent}{0pt}
    \setlength{\parskip}{6pt plus 2pt minus 1pt}}
}{% if KOMA class
  \KOMAoptions{parskip=half}}
\makeatother
\usepackage{xcolor}
\usepackage{color}
\usepackage{fancyvrb}
\newcommand{\VerbBar}{|}
\newcommand{\VERB}{\Verb[commandchars=\\\{\}]}
\DefineVerbatimEnvironment{Highlighting}{Verbatim}{commandchars=\\\{\}}
% Add ',fontsize=\small' for more characters per line
\newenvironment{Shaded}{}{}
\newcommand{\AlertTok}[1]{\textcolor[rgb]{1.00,0.00,0.00}{\textbf{#1}}}
\newcommand{\AnnotationTok}[1]{\textcolor[rgb]{0.38,0.63,0.69}{\textbf{\textit{#1}}}}
\newcommand{\AttributeTok}[1]{\textcolor[rgb]{0.49,0.56,0.16}{#1}}
\newcommand{\BaseNTok}[1]{\textcolor[rgb]{0.25,0.63,0.44}{#1}}
\newcommand{\BuiltInTok}[1]{\textcolor[rgb]{0.00,0.50,0.00}{#1}}
\newcommand{\CharTok}[1]{\textcolor[rgb]{0.25,0.44,0.63}{#1}}
\newcommand{\CommentTok}[1]{\textcolor[rgb]{0.38,0.63,0.69}{\textit{#1}}}
\newcommand{\CommentVarTok}[1]{\textcolor[rgb]{0.38,0.63,0.69}{\textbf{\textit{#1}}}}
\newcommand{\ConstantTok}[1]{\textcolor[rgb]{0.53,0.00,0.00}{#1}}
\newcommand{\ControlFlowTok}[1]{\textcolor[rgb]{0.00,0.44,0.13}{\textbf{#1}}}
\newcommand{\DataTypeTok}[1]{\textcolor[rgb]{0.56,0.13,0.00}{#1}}
\newcommand{\DecValTok}[1]{\textcolor[rgb]{0.25,0.63,0.44}{#1}}
\newcommand{\DocumentationTok}[1]{\textcolor[rgb]{0.73,0.13,0.13}{\textit{#1}}}
\newcommand{\ErrorTok}[1]{\textcolor[rgb]{1.00,0.00,0.00}{\textbf{#1}}}
\newcommand{\ExtensionTok}[1]{#1}
\newcommand{\FloatTok}[1]{\textcolor[rgb]{0.25,0.63,0.44}{#1}}
\newcommand{\FunctionTok}[1]{\textcolor[rgb]{0.02,0.16,0.49}{#1}}
\newcommand{\ImportTok}[1]{\textcolor[rgb]{0.00,0.50,0.00}{\textbf{#1}}}
\newcommand{\InformationTok}[1]{\textcolor[rgb]{0.38,0.63,0.69}{\textbf{\textit{#1}}}}
\newcommand{\KeywordTok}[1]{\textcolor[rgb]{0.00,0.44,0.13}{\textbf{#1}}}
\newcommand{\NormalTok}[1]{#1}
\newcommand{\OperatorTok}[1]{\textcolor[rgb]{0.40,0.40,0.40}{#1}}
\newcommand{\OtherTok}[1]{\textcolor[rgb]{0.00,0.44,0.13}{#1}}
\newcommand{\PreprocessorTok}[1]{\textcolor[rgb]{0.74,0.48,0.00}{#1}}
\newcommand{\RegionMarkerTok}[1]{#1}
\newcommand{\SpecialCharTok}[1]{\textcolor[rgb]{0.25,0.44,0.63}{#1}}
\newcommand{\SpecialStringTok}[1]{\textcolor[rgb]{0.73,0.40,0.53}{#1}}
\newcommand{\StringTok}[1]{\textcolor[rgb]{0.25,0.44,0.63}{#1}}
\newcommand{\VariableTok}[1]{\textcolor[rgb]{0.10,0.09,0.49}{#1}}
\newcommand{\VerbatimStringTok}[1]{\textcolor[rgb]{0.25,0.44,0.63}{#1}}
\newcommand{\WarningTok}[1]{\textcolor[rgb]{0.38,0.63,0.69}{\textbf{\textit{#1}}}}
\usepackage{longtable,booktabs,array}
\usepackage{calc} % for calculating minipage widths
% Correct order of tables after \paragraph or \subparagraph
\usepackage{etoolbox}
\makeatletter
\patchcmd\longtable{\par}{\if@noskipsec\mbox{}\fi\par}{}{}
\makeatother
% Allow footnotes in longtable head/foot
\IfFileExists{footnotehyper.sty}{\usepackage{footnotehyper}}{\usepackage{footnote}}
\makesavenoteenv{longtable}
\setlength{\emergencystretch}{3em} % prevent overfull lines
\providecommand{\tightlist}{%
  \setlength{\itemsep}{0pt}\setlength{\parskip}{0pt}}
\setcounter{secnumdepth}{5}
\usepackage[margin=2.1cm]{geometry}
\usepackage{longtable,booktabs,array}
\usepackage{microtype}
\setlength{\emergencystretch}{4em}
\sloppy
\ifLuaTeX
  \usepackage{selnolig}  % disable illegal ligatures
\fi
\usepackage{bookmark}
\IfFileExists{xurl.sty}{\usepackage{xurl}}{} % add URL line breaks if available
\urlstyle{same}
\hypersetup{
  hidelinks,
  pdfcreator={LaTeX via pandoc}}

\author{}
\date{}

\begin{document}

{
\setcounter{tocdepth}{3}
\tableofcontents
}
\section{DDTA - Working Metamodel - Functional
Requirement}\label{ddta---working-metamodel---functional-requirement}

\textbf{NON-CANONICAL WORKING DRAFT - REVISION 1}

\subsection{1. Status and purpose}\label{status-and-purpose}

This document starts the Functional Requirement (FR) semantic study
after the Macro Requirement and Decision models reached a provisional
upper-layer baseline.

The purpose is not to reproduce ThreatForge's current Requirement
schema. It is to identify the smallest conceptual model that explains
what a Functional Requirement means, how it relates vertically to Macro
Requirements and Decisions, and how it connects downstream to
realization and verification without embedding implementation into
requirement semantics.

For this phase, \textbf{Requirement means Functional Requirement unless
explicitly qualified}. Governance, Security and Privacy are deferred
specialized-requirement studies.

The current upper-layer baseline remains reopenable. A concrete FR
counterexample may require a correction to MR or Decision semantics or
relations, followed by regression over prior evidence.

\subsection{2. Evidence boundary}\label{evidence-boundary}

\subsubsection{2.1 Admitted construction
evidence}\label{admitted-construction-evidence}

ThreatForge product-semantic evidence is read at commit:

\texttt{cae0f7b6b37f430ac4e857aabf6ef9f87c89dbb1}

Construction evidence is limited to ThreatForge MR-0001 and MR-0002 plus
the already admitted independent facial-access MR/Decision example.

Selected historical sources used to seed questions, not to define truth,
are the active Functional Requirement model/profile, the MR-0001/MR-0002
Requirement registries, and two representative bodies. Their filenames
are:

\begin{itemize}
\tightlist
\item
  \texttt{functional-requirement.model.yml}
\item
  \texttt{functional-requirement-body.profile.yml}
\item
  \texttt{MR-0001.requirements.registry.yml}
\item
  \texttt{MR-0002.requirements.registry.yml}
\item
  \texttt{MR-0001ADR-0007REQ-0001\_body.md}
\item
  \texttt{MR-0002ADR-0003REQ-0001\_body.md}
\end{itemize}

All are under the canonical \texttt{docs/reference/project-model/} tree
at the pinned ThreatForge commit above.

ThreatForge currently describes the Functional Requirement model as an
``independently testable governed obligation'' and gives every FR
registry record both a Macro Requirement ID and Decision ID. Its
Markdown profile requires \texttt{Intent},
\texttt{Functional\ obligation}, \texttt{Scope} and \texttt{Acceptance}
sections. These are historical design choices to test, not automatically
adopted L1 semantics.

\subsubsection{2.2 Holdout protection}\label{holdout-protection}

Do not inspect Requirement content belonging to ThreatForge MR-0003,
MR-0004 or MR-0005 while constructing this model. Their unconsumed
Decision/Requirement families remain future vertical holdout evidence.

\subsection{3. Fundamental question}\label{fundamental-question}

The working FR question is:

\begin{quote}
\textbf{What must be observably or objectively true of the governed
project/system so that a macro concern and its applicable choices become
operational, without prescribing the concrete implementation that makes
it true?}
\end{quote}

This question deliberately separates:

\begin{Shaded}
\begin{Highlighting}[]
\NormalTok{MacroRequirement}
\NormalTok{    major concern / intended result}

\NormalTok{Decision}
\NormalTok{    significant choice that narrows the concern}

\NormalTok{FunctionalRequirement}
\NormalTok{    functional truth/capability/behavior/result that must hold}

\NormalTok{Realization}
\NormalTok{    concrete artifacts and mechanisms that make it hold}

\NormalTok{Verification}
\NormalTok{    evidence that supports the claim that it holds}
\end{Highlighting}
\end{Shaded}

\subsection{4. Working semantic
definition}\label{working-semantic-definition}

\textbf{FR-DEF-01 - Functional Requirement, working definition}

A Functional Requirement is a governed, independently assessable
obligation stating one coherent capability, behavior, service, state
transition or result that the governed project/system must provide or
make true. It may operationalize project intent and applicable
Decisions, but it does not prescribe the concrete realization mechanism
unless that mechanism is itself intrinsic to the required
externally/operationally meaningful function.

This definition is a construction hypothesis, not closure.

\subsubsection{4.1 Why ``independently
assessable''}\label{why-independently-assessable}

A Requirement should permit a meaningful determination of satisfaction
that is stronger than ``some implementation artifact exists''.
Assessment may ultimately use tests, checks, inspection or other
governed evidence, but the exact verification-evidence model is
downstream.

\subsubsection{4.2 Why ``one coherent''
function}\label{why-one-coherent-function}

A Requirement may need multiple normative statements to define one
capability. Multiple sentences do not automatically imply multiple
Requirements.

Use the split test:

\begin{quote}
If two obligation groups can be realized, accepted or changed
independently and neither is necessary to define the same indivisible
functional capability, review them as separate Functional Requirements.
\end{quote}

This mirrors the Decision coherence test without assuming one bullet
equals one semantic entity.

\subsection{5. Candidate semantic shape}\label{candidate-semantic-shape}

Revision 1 intentionally does \textbf{not} freeze the historical
ThreatForge body sections as the FR semantic core.

\begin{Shaded}
\begin{Highlighting}[]
\NormalTok{FunctionalRequirement}
\NormalTok{|{-} title?                 [likely common governed{-}document concern]}
\NormalTok{|{-} functionalObligation  [strong L1 candidate]}
\NormalTok{|{-} scope?                 [candidate; requires evidence]}
\NormalTok{\textasciigrave{}{-} acceptance?            [candidate; may belong to verification contract]}
\end{Highlighting}
\end{Shaded}

\texttt{Intent}, present in the historical ThreatForge representation,
is not yet admitted to the minimum semantic core. Its value must be
tested against information already recoverable from the parent context
plus the functional obligation.

\subsubsection{5.1 Candidate-field status}\label{candidate-field-status}

\begin{longtable}[]{@{}
  >{\raggedright\arraybackslash}p{(\columnwidth - 4\tabcolsep) * \real{0.3333}}
  >{\raggedright\arraybackslash}p{(\columnwidth - 4\tabcolsep) * \real{0.3333}}
  >{\raggedright\arraybackslash}p{(\columnwidth - 4\tabcolsep) * \real{0.3333}}@{}}
\toprule\noalign{}
\begin{minipage}[b]{\linewidth}\raggedright
Candidate
\end{minipage} & \begin{minipage}[b]{\linewidth}\raggedright
Revision 1 status
\end{minipage} & \begin{minipage}[b]{\linewidth}\raggedright
Question
\end{minipage} \\
\midrule\noalign{}
\endhead
\bottomrule\noalign{}
\endlastfoot
\texttt{title} & common identity/display candidate & Is title intrinsic
FR semantics or common \texttt{GovernedDocument} semantics? \\
\texttt{functionalObligation} & \textbf{strong required candidate} & Can
any valid FR exist without stating what must be true/provided? \\
\texttt{scope} & open & Does scope carry independent semantic
information, or can it be expressed by the obligation/context without
repetition? \\
\texttt{acceptance} & open & Is acceptance intrinsic Requirement
semantics or a separate verification contract/evidence relation? \\
\texttt{intent} & open / suspected non-core & Does FR-local intent add
non-redundant meaning after MR + Decision + obligation are known? \\
\end{longtable}

No representation ordering, Markdown heading, normative keyword or YAML
member is L1 merely because ThreatForge currently enforces it.

\subsection{6. The central vertical relation is still
open}\label{the-central-vertical-relation-is-still-open}

Historical ThreatForge encodes:

\begin{Shaded}
\begin{Highlighting}[]
\NormalTok{MR {-}\textgreater{} Decision {-}\textgreater{} FunctionalRequirement}
\end{Highlighting}
\end{Shaded}

with one \texttt{decision\_id} and one \texttt{macro\_requirement\_id}
in each FR registry record. The current logical model treats both as
relations.

Revision 1 does \textbf{not} close that topology yet.

Two competing hypotheses must be tested.

\subsubsection{H-FR-PARENT-A - Mandatory Decision
parent}\label{h-fr-parent-a---mandatory-decision-parent}

\begin{Shaded}
\begin{Highlighting}[]
\NormalTok{MacroRequirement}
\NormalTok{    1}
\NormalTok{    |}
\NormalTok{    *}
\NormalTok{Decision}
\NormalTok{    1}
\NormalTok{    |}
\NormalTok{    *}
\NormalTok{FunctionalRequirement}
\end{Highlighting}
\end{Shaded}

Every FR is the operational downstream expression of exactly one
significant governed choice.

Potential advantage: strong vertical trace and simple ownership.

Potential defect: it may force artificial Decisions for functions that
follow directly from the Macro Requirement/domain obligation without an
unresolved choice.

\subsubsection{H-FR-PARENT-B - MR ownership plus Decision
applicability}\label{h-fr-parent-b---mr-ownership-plus-decision-applicability}

\begin{Shaded}
\begin{Highlighting}[]
\NormalTok{MacroRequirement 1 {-}{-}{-}{-} * FunctionalRequirement}

\NormalTok{Decision * {-}{-}{-}{-} * FunctionalRequirement}
\NormalTok{          constrains / is operationalized by}
\end{Highlighting}
\end{Shaded}

Every FR has one macro-concern owner, while zero or more Decisions may
constrain or be operationalized by it.

Potential advantage: does not require dummy Decisions.

Potential defect: may weaken the simple descending hierarchy and
requires an explicit model for applicable Decision context.

The construction and independent examples must decide between these
hypotheses or expose a smaller third relation model.

\subsubsection{6.1 Direct MR identifier may be
derived}\label{direct-mr-identifier-may-be-derived}

Even if H-FR-PARENT-A survives, a stored \texttt{macro\_requirement\_id}
may be a derived/indexed mirror of the Decision parent rather than a
second independent semantic parent relation.

Identity serialization must not be used as proof of topology.

\subsection{7. Candidate FR invariants}\label{candidate-fr-invariants}

These are working tests. Only \texttt{FR-CAND-01} through
\texttt{FR-CAND-05} are strong enough to guide initial authoring; none
is closed before regression.

\subsubsection{FR-CAND-01 - Functional
obligation}\label{fr-cand-01---functional-obligation}

An FR states a required capability, behavior, service, transition or
result. Pure rationale, architecture preference, implementation task or
evidence record is not an FR.

\subsubsection{FR-CAND-02 - Independent
assessability}\label{fr-cand-02---independent-assessability}

There must be a meaningful way to decide whether the required functional
truth holds. ``Code exists'' or ``file created'' is not sufficient
unless file creation itself is the externally/operationally required
behavior.

\subsubsection{FR-CAND-03 - Coherent requirement
unit}\label{fr-cand-03---coherent-requirement-unit}

The FR owns one coherent functional capability. Split independently
realizable/acceptable capability families.

\subsubsection{FR-CAND-04 - Decision
non-repetition}\label{fr-cand-04---decision-non-repetition}

An FR does not merely restate the parent/applicable Decision. The
Decision explains which choice was made; the FR states what must
consequently function or be true.

\subsubsection{FR-CAND-05 - Implementation
independence}\label{fr-cand-05---implementation-independence}

The FR does not prescribe frameworks, source paths, internal
class/function layout, code annotations, transport technology or other
realization detail unless that mechanism is itself part of the required
observable contract.

\subsubsection{FR-CAND-06 - Realization is a relation, not FR
content}\label{fr-cand-06---realization-is-a-relation-not-fr-content}

Implementation artifacts may realize an FR, but their existence, paths
and implementation state belong to realization trace/governance rather
than the intrinsic FR meaning.

\subsubsection{FR-CAND-07 - Verification is distinct from
realization}\label{fr-cand-07---verification-is-distinct-from-realization}

Evidence supporting satisfaction is not the same thing as an artifact
that realizes the Requirement.

\subsubsection{FR-CAND-08 - Specialized constraints do not replace the
FR}\label{fr-cand-08---specialized-constraints-do-not-replace-the-fr}

Security, Governance and Privacy Requirements may enrich/constrain an
FR. The underlying functional capability remains owned by the FR.

\subsubsection{FR-CAND-09 - Effective obligation is derived
context}\label{fr-cand-09---effective-obligation-is-derived-context}

For realization/review, the effective obligation of an FR includes the
FR plus all applicable specialized Requirements. This is a derived
governed view, not a reason to copy specialized constraints into the FR
body.

\subsubsection{FR-CAND-10 - Corpus non-redundancy
candidate}\label{fr-cand-10---corpus-non-redundancy-candidate}

Two FRs that express the same functional truth under the same effective
governed context may indicate duplication or incorrect decomposition. A
corpus-level invariant analogous to Decision \texttt{H-CLOSED-02} is
plausible but remains open until multiple FR families are analyzed.

\subsection{8. Realization and verification
boundary}\label{realization-and-verification-boundary}

A Functional Requirement should lead toward implementation without
becoming implementation.

Conceptually:

\begin{Shaded}
\begin{Highlighting}[]
\NormalTok{FunctionalRequirement}
\NormalTok{      |}
\NormalTok{      | realizedBy}
\NormalTok{      v}
\NormalTok{RealizationArtifact *}
\end{Highlighting}
\end{Shaded}

and independently:

\begin{Shaded}
\begin{Highlighting}[]
\NormalTok{Requirement}
\NormalTok{      |}
\NormalTok{      | verifiedBy}
\NormalTok{      v}
\NormalTok{VerificationEvidence *}
\end{Highlighting}
\end{Shaded}

Both cardinalities are deliberately open at L1 because an authored
Requirement may exist before implementation/evidence. Lifecycle may
later impose stronger conditions for states such as
implemented/accepted.

One implementation artifact may realize several Requirements, and one
Requirement may require several artifacts. Therefore no one-file-per-FR
assumption is admitted.

\subsubsection{8.1 RealizationArtifact is intentionally
generic}\label{realizationartifact-is-intentionally-generic}

The portable metamodel should not assume that realization means source
code only. A functional capability may be realized by code,
configuration, infrastructure, a composed service, or another project
artifact.

ThreatForge may later support specific artifact classes without making
those classes universal FR semantics.

\subsection{9. Specialized Requirement
boundary}\label{specialized-requirement-boundary}

Revision 1 reserves the following future structure without closing its
common semantics:

\begin{Shaded}
\begin{Highlighting}[]
\NormalTok{FunctionalRequirement}
\NormalTok{      |}
\NormalTok{      +{-}{-} GovernanceRequirement *}
\NormalTok{      +{-}{-} SecurityRequirement *}
\NormalTok{      \textasciigrave{}{-}{-} PrivacyRequirement *}
\end{Highlighting}
\end{Shaded}

The current working interpretation is:

\begin{quote}
A Specialized Requirement adds an obligation/constraint that the
realization and/or verification of its Functional Requirement must
satisfy, without replacing the functional capability owned by the FR.
\end{quote}

Whether \texttt{specializes}, \texttt{constrains}, \texttt{appliesTo} or
another relation name best represents this meaning remains open.

\subsubsection{9.1 Security-specific future
origin}\label{security-specific-future-origin}

Security is not created merely by the FR hierarchy. The later Security
model must also preserve analytical provenance:

\begin{Shaded}
\begin{Highlighting}[]
\NormalTok{AnalysisRecord}
\NormalTok{    |}
\NormalTok{    v}
\NormalTok{CommonFinding}
\NormalTok{    |}
\NormalTok{    | accepted}
\NormalTok{    v}
\NormalTok{SecurityRequirement}
\NormalTok{    |}
\NormalTok{    v}
\NormalTok{FunctionalRequirement}
\end{Highlighting}
\end{Shaded}

The Security Requirement states a product/project obligation. STRIDE or
STRIDE-AI remains recoverable through the Analysis Record rather than
becoming a field of the FR.

\subsection{10. Independent facial-access
probe}\label{independent-facial-access-probe}

The facial-access example deliberately creates a vertical probe that may
challenge mandatory Decision parentage.

\subsubsection{10.1 Existing upper-layer
context}\label{existing-upper-layer-context}

Macro concern (facial-access example): facial recognition should
determine whether the acquired face corresponds to the expected
authorized person with sufficient quality for the access process.

Existing admitted Decision proposition from the Decision study: the
recognition result consumed by the access process has stable governed
semantics independent from raw/model-specific ML score semantics.

\subsubsection{10.2 FR proposition A - acquisition suitable for
verification}\label{fr-proposition-a---acquisition-suitable-for-verification}

\textbf{Working obligation}

\begin{quote}
The system must acquire facial information at the access point that is
sufficient to attempt person verification and must make an insufficient
acquisition distinguishable from a valid verification result.
\end{quote}

This looks locally like a Functional Requirement: it states a functional
capability and can be assessed without naming a camera model, library or
source file.

But its parentage is intentionally unresolved. It may follow directly
from the Macro Requirement rather than from the specific Decision about
stable result semantics.

Therefore this FR is a direct test of H-FR-PARENT-A versus
H-FR-PARENT-B.

\subsubsection{10.3 FR proposition B - governed verification
result}\label{fr-proposition-b---governed-verification-result}

\textbf{Working obligation}

\begin{quote}
For every completed verification attempt, the system must expose a
result using the governed verification semantics consumed by the access
process rather than exposing raw model-specific score semantics as the
functional result.
\end{quote}

This proposition appears to operationalize the existing stable-semantics
Decision much more directly.

The pair is useful because both are legitimate functions under the same
macro concern, while only one is obviously downstream of the known
Decision.

\subsubsection{10.4 Implementation is
separate}\label{implementation-is-separate}

Possible realization artifacts such as a camera adapter, capture
service, face-verification service or score interpreter are not part of
either FR definition.

Likewise, tests for insufficient capture or stable result mapping are
verification evidence, not the FR itself.

\subsection{11. ThreatForge construction
probes}\label{threatforge-construction-probes}

\subsubsection{11.1 Scaffold Requirement}\label{scaffold-requirement}

Historical \texttt{MR-0002ADR-0003REQ-0001} states a coherent user/tool
capability - create a governed implementation scaffold - but its current
Functional obligation also includes details such as specific source
declarations, a deterministic artifact identifier, implementation-trace
mutation and invocation through a VS Code task.

The FR study must classify these statements rather than assume the whole
historical body is one semantic Requirement:

\begin{itemize}
\tightlist
\item
  required functional behavior;
\item
  consequence of an architectural Decision;
\item
  realization/trace mechanism;
\item
  specialized Governance constraint;
\item
  separate FR that can change independently.
\end{itemize}

This is a strong construction case for \texttt{FR-CAND-03},
\texttt{FR-CAND-04} and \texttt{FR-CAND-05}.

\subsubsection{11.2 Canonical profiles
Requirement}\label{canonical-profiles-requirement}

Historical \texttt{MR-0001ADR-0007REQ-0001} requires canonical logical
models and representation profiles for active governed document models.
It is a useful contrast because its obligation is closer to a
product/tool capability than to a concrete implementation path, yet
several bullets may still represent separately changeable obligations.

It therefore probes the semantic granularity of one FR versus a family
of FRs.

\subsection{12. Authoring diagnostics for the construction
phase}\label{authoring-diagnostics-for-the-construction-phase}

Use these questions when rewriting or authoring candidate FRs:

\begin{enumerate}
\def\labelenumi{\arabic{enumi}.}
\tightlist
\item
  \textbf{Functional truth:} What must actually work or be true?
\item
  \textbf{Choice test:} Is this sentence still a choice/trade-off? If
  yes, it may belong to Decision.
\item
  \textbf{Implementation test:} Could a different implementation satisfy
  the same sentence? If not, is the named mechanism genuinely part of
  the required contract or merely realization detail?
\item
  \textbf{Split test:} Could one obligation change/be accepted without
  the other? If yes, consider separate FRs.
\item
  \textbf{Parent test:} Which MR/Decision context naturally owns this
  function? Would forcing one Decision parent create a dummy Decision?
\item
  \textbf{Verification test:} What observation could establish
  satisfaction without treating artifact existence as proof?
\item
  \textbf{Specialization test:} Is this the function itself, or an
  additional Security/Governance/Privacy constraint on that function?
\item
  \textbf{Copy test:} Is the FR repeating a Decision, controlled
  vocabulary or specialized constraint that should be
  referenced/effectively inherited instead?
\end{enumerate}

\subsection{13. What Revision 1 deliberately does not
decide}\label{what-revision-1-deliberately-does-not-decide}

\begin{itemize}
\tightlist
\item
  whether every FR has exactly one Decision parent;
\item
  whether FR has an independent semantic MR relation or only a
  derived/mirrored one;
\item
  whether \texttt{Intent} is an FR field;
\item
  whether \texttt{Scope} is semantically required;
\item
  whether acceptance conditions are intrinsic FR semantics or a separate
  verification contract;
\item
  the exact relation names/cardinalities for realization and
  verification;
\item
  FR lifecycle/status;
\item
  FR-to-FR dependencies;
\item
  the common superclass semantics of Governance/Security/Privacy;
\item
  Privacy-specific origin semantics;
\item
  Security-specific Finding cardinality beyond the already observed
  future derivation requirement;
\item
  any YAML/Markdown representation profile;
\item
  any ThreatForge implementation mechanism.
\end{itemize}

\subsection{14. Construction sequence from this
draft}\label{construction-sequence-from-this-draft}

\begin{Shaded}
\begin{Highlighting}[]
\NormalTok{Revision 1 FR hypotheses}
\NormalTok{    |}
\NormalTok{    v}
\NormalTok{inventory every Functional Requirement in ThreatForge MR{-}0001/MR{-}0002}
\NormalTok{    |}
\NormalTok{    v}
\NormalTok{classify each statement:}
\NormalTok{FR semantics / Decision / realization / verification / specialization / representation}
\NormalTok{    |}
\NormalTok{    v}
\NormalTok{test parent hypothesis and requirement granularity}
\NormalTok{    |}
\NormalTok{    v}
\NormalTok{author independent facial{-}access FR corpus}
\NormalTok{    |}
\NormalTok{    v}
\NormalTok{vertical MR {-}\textgreater{} Decision {-}\textgreater{} FR regression}
\NormalTok{    |}
\NormalTok{    +{-}{-}\textgreater{} no counterexample: strengthen/close candidates}
\NormalTok{    |}
\NormalTok{    \textasciigrave{}{-}{-}\textgreater{} counterexample: correct smallest owning layer and regress prior evidence}
\end{Highlighting}
\end{Shaded}

Only after the Functional Requirement boundary is stable should the
study define the shared \texttt{SpecializedRequirement} semantics and
then test Security/Governance/Privacy individually.

\subsection{15. Revision 1 working
summary}\label{revision-1-working-summary}

The smallest current hypothesis is intentionally modest:

\begin{Shaded}
\begin{Highlighting}[]
\NormalTok{FunctionalRequirement}
\NormalTok{    = governed functional obligation}
\NormalTok{    + independently assessable satisfaction}
\NormalTok{    + no intrinsic implementation prescription}
\end{Highlighting}
\end{Shaded}

The major unresolved structural question is whether the vertical
topology is truly:

\begin{Shaded}
\begin{Highlighting}[]
\NormalTok{MR {-}\textgreater{} Decision {-}\textgreater{} FR}
\end{Highlighting}
\end{Shaded}

or whether the evidence requires a weaker relation in which MR owns the
FR and Decisions constrain/are operationalized by it.

That question should be answered by evidence before the existing
ThreatForge ID grammar or registry relations are promoted into the DDTA
conceptual metamodel.

\end{document}
